% Der Hauptsatz gehört zum Fachband, die sechs Hilfsresultate zum Beweisband.
% Abschnittsnummern und anonyme Anker der späteren Ergebnisse bleiben stabil.
\section{Produkte endlicher Wörter}

Für eine Halbgruppe \(\SemigroupAlg{A}{\star}\) bezeichnet
\(\WordFold{\star}{u}\) die in Band~27 rekursiv definierte Linksfaltung
des nichtleeren Wortes \(u\in\NonemptyWords{A}\). Sie wertet die Faktoren
von links nach rechts aus. Das
\BracketingResultLink{SemigroupWordBlockLaw}{Blockgesetz in den Beweistabellen}
zeigt, dass man auch zunächst zusammenhängende Blöcke auswerten und danach
deren Werte verknüpfen darf.

% Die drei bisherigen Satzanker reservieren, ohne Aussagen zu deklarieren.
\phantomsection\phantomsection\phantomsection

\section{Klammerungsunabhängigkeit}

Alle vollständigen Klammerungen desselben nichtleeren Wortes liefern in
einer Halbgruppe denselben Wert. Der folgende Hauptsatz hält diese Aussage
in der Sprache der Klammerungsbäume fest.

\Needspace{16\baselineskip}
% Drei ausgelagerte Hilfssatzanker; der Hauptsatz behält seine Originalnummer.
\phantomsection\phantomsection\phantomsection
\setcounter{formulaThm}{3}
\hypertarget{bracketing.main.independence}{}
\BracketingIndependence

\noindent\emph{Beweis:} Die \BracketingReadingLink{} erklärt die
Beweiskette mit Beispielen, Diagrammen und den beiden Induktionsargumenten.
Die \BracketingProofsLink{} enthalten die sechs Hilfsresultate und alle
Ableitungen einschließlich des \BracketingIndependenceProofLink{}.

\par\medskip
\hypertarget{bracketing.notation}{}

\noindent\textbf{Anmerkung zur Produktschreibweise.}
In einer Halbgruppe haben alle vollständigen Klammerungen eines endlichen,
nichtleeren Produkts denselben Wert, sofern die Reihenfolge der Faktoren
unverändert bleibt. Für jede natürliche Zahl \(n\geq 1\) und Faktoren
\(a_1,\ldots,a_n\in A\) dürfen wir deshalb einfach
\[
  a_1\star a_2\star\cdots\star a_n
\]
schreiben. Gemeint ist der gemeinsame Wert aller dieser Klammerungen,
beispielsweise der von links ausgewertete Wert. Bei einem einzigen Faktor
ist das Produkt gleich \(a_1\).

Die Klammern können somit weggelassen oder nach Bedarf gesetzt werden.
Ein beliebiges Vertauschen von Faktoren ist damit nicht erlaubt; dazu wäre zusätzlich
Kommutativität nötig. Ein leeres Produkt wird hier nicht eingeführt, da
eine Halbgruppe kein neutrales Element besitzen muss.

Die im Beweis verwendete
\BracketingResultLink{SemigroupTreeNormalForm}{Baum-Normalform}, die jede
Baumauswertung auf die Linksfaltung ihres Blattwortes zurückführt,
steht zusammen mit ihrem Beweis in den Beweistabellen.
